% ===================== RESUMEN EJECUTIVO =====================
\newpage
\section*{Resumen Ejecutivo}
\addcontentsline{toc}{section}{Resumen Ejecutivo}

El presente informe recopila la investigación técnica y de mercado ejecutada para la selección de los equipos del sistema de ventilación y presurización positiva del laboratorio de análisis industrial de BRINSA, ubicado en Cajicá, Cundinamarca, a 2 558 msnm, dentro de una planta de hipoclorito de calcio de ambiente exterior altamente corrosivo. El sistema impulsa 3 840 m$^3$/h de aire filtrado (12 renovaciones/h sobre un volumen efectivo de 320 m$^3$) y mantiene el recinto a +25 Pa respecto al exterior, con mínimo admisible de +12.5 Pa, para excluir insectos, objetos extraños y polvo; al no ser un sistema de biocontención, la filtración HEPA se descarta y la etapa definitiva es MERV 13-14.

El resultado técnico central de la investigación es la corrección de densidad por altitud: la presión de estación real es 74.1 kPa y la densidad del aire 0.88 kg/m$^3$, lo que fija un factor de densidad $k$ = 0.733 respecto a las curvas de catálogo. En consecuencia, el ventilador se selecciona en el punto equivalente de catálogo 3 840 m$^3$/h a 260 Pa (190 Pa en el sitio, escenario de filtro cargado), con una potencia teórica de aire de 0.338 kW y un motor instalado de 1.0 HP TEFC anticorrosivo, 440 V, 3$\phi$, 60 Hz. La misma corrección reduce la presión sostenida naturalmente por las rejillas a 11 Pa, por lo que el damper de alivio barométrico calibrable pasa de opción a componente obligatorio del sistema.

La selección recomendada comprende: ventilador centrífugo de álabes curvados hacia atrás en PRFV (Greenheck BCSW-FRP como primera opción, con canal local en Bogotá); filtración en dos etapas, prefiltro MERV 8 y filtro final V-bank MERV 13-14 de marco plástico (Camfil Durafil ES2 como referencia de diseño); tres rejillas eggcrate de 353 mm $\times$ 336 mm en acero inoxidable 316L con malla anti-insectos inox 316 de 18$\times$18; damper barométrico Greenheck SEBR-10 calibrado a +25 Pa en comisionamiento; y transmisor Dwyer MS-121 (0-62.5 Pa, 4-20 mA, NEMA 4X) con alarmas a +12.5 Pa y +40 Pa, complementado con un Magnehelic de lectura local. Todas las fuentes comerciales y técnicas se citan con URL y fecha de consulta (23 de julio de 2026).

\newpage
